% This section lists a number of 2-3 research questions that the thesis should address. 
% The research questions should be put into the context of the idea, and you should explain 
% what they mean and how you want to address each. We also suggest including an example scenario 
% of a successful outcome of the thesis, i.e., what would your tool allow us to do once it exists?

\chapter{Research questions} \label{research_questions}


This paper will address the following research questions:
\begin{enumerate}
    \item How can automated model checking and other formal-verification 
        techniques be used for verifying liveness (replication reliability),
        safety and other properties in self-replicating robotic systems(SRRS), 
        and what domain-specific adaptations are required for practical implementation?
    \item What systematic methodology can translate the behaviors and 
        structural constraints of self-replicating robots into expressive 
        temporal-logic models and specifications?
    \item What limitations—computational, semantic, or 
        modelling-related -- does developed approach face, and how might these limitations be mitigated?\\
        Compare the  finite state machine algorithm control with SMV counterexample-based control 
        in simulation by time, flexibility

\end{enumerate}


\section{Analysis of model checking approaches in self-replicating systems domain}

% The first question includes examination of the state-of-the-art 
% of model checking approaches to determine whether they can be used to model and verify 
% an SRR system.
%
% The experience and different practical approaches in the field of self-reproducing 
% automata and robots should be studied, based on which the main 
% properties that should be included in the specifications have to be identified.
%
% Initially can be considered basic properties like liveness (replication reliability) and safety.
% With some adaptations it should be possible to express it using every model-checking
% specification language.
% During this analysis or further experimental checking, 
% more complex properties can be discovered  and integrated.
The first research question involves a comprehensive survey and critical 
evaluation of current state-of-the-art model checking and formal verification 
tools, such as SPIN, NuSMV, and PRISM, focusing on their applicability to modeling 
and verifying the unique dynamics of self-replicating robotic systems (SRRS). 

This investigation includes a review of relevant literature and existing systems 
in self-replicating robotics and cellular automata to identify key operational 
characteristics and formalizable behaviors. It also involves identifying 
core properties relevant to SRRS, such as liveness (ensuring successful 
replication within time or resource constraints), safety (avoiding 
resource conflicts or hazardous outcomes), as well as broader system-level guarantees 
like termination, deadlock-freedom, and fault-tolerance. 
The study will further explore how these properties can be 
expressed in various formal languages, including Linear 
Temporal Logic (LTL), Computation Tree Logic (CTL), and probabilistic extensions. 
This will require examining what adaptations or abstractions are needed to 
make these languages suitable for SRRS modeling. 

Iterative experimentation with selected tools will help assess the feasibility 
of encoding and verifying these properties and may lead to the discovery 
of more complex or emergent behaviors to integrate in later stages.

\section{Experimental checking}



The second research question aims to define a systematic formal modeling 
and verification workflow tailored to SRRS. 
This involves developing a formal representation of 
self-replication processes, structural rules, environmental interactions, 
and resource usage. 

It also includes establishing a methodology to map real-world 
or simulated robotic behaviors - such as perception, manipulation, and actuation- onto 
abstract states and transitions that are amenable to formal analysis. 
Based on this framework, appropriate model-checking tools and specification 
languages will be selected and justified through experimental evaluation within 
simulation environments. 

A prototype or proof-of-concept implementation will be developed 
using a simplified SRRS scenario, such modular robotic replication, 
to demonstrate how this methodology captures essential properties and enables formal verification. 

The ultimate goal is to bridge the gap between low-level robotic simulation and control 
systems and high-level formal models, ensuring that the designed 
behaviors are both physically feasible and logically justified.


\section{Limitations and challenges}


The third research question addresses the fundamental limitations that 
current formal methods and model-checking tools face when applied to the 
complex domain of self-replicating robotic systems.
It includes analyzing computational constraints such as state-space explosion, 
undecidability in dynamic conditions, and potential inefficiencies in scaling existing verification tools. 

The study will also examine semantic mismatches between physical 
or simulated robot behavior and abstract models, particularly in representing kinematic 
and dynamic aspects, sensor errors, and asynchronous or partially 
observable environments. 

Additionally, the research will explore modeling limitations related 
to the representation of dynamic structure generation and adaptive behaviors.

To address these challenges, various mitigation strategies will be proposed,
including modular verification, changes in system abstractions  
and hybrid techniques that combine formal models with simulation-based validation. 

The outcome of this part may include recommendations for extending the 
selected verification toolsets, adapting implementation strategies, or designing 
new experimental setups specifically for the analysis of SRRS.


% This section lists a number of 2-3 research questions that the thesis should address. 
% The research questions should be put into the context of the idea, and you should explain 
% what they mean and how you want to address each. We also suggest including an example scenario 
% of a successful outcome of the thesis, i.e., what would your tool allow us to do once it exists?

\chapter{Methodology}




\section{Research tools}


Self-replication research in robotics involves complex interactions between 
assembly processes, modular design, environmental constraints, 
and control systems. 
Robotic simulation environments offer a valuable sandbox to 
explore these dynamics before physical implementation. 
The ideal simulation platform must support modularity, kinematics, 
sensing, and interaction with discrete components.

A set of tools for simulation and control of robotic systems including SRRS can be 
seen in the Table \ref{tab:tools}

\begin{table}[h!]
    \centering
    \begin{tabular}{|m{2cm}|m{1.5cm}|m{2cm}|m{2.5cm}|m{5cm}|}
        \hline
        \textbf{Simulator} & \textbf{Physics fidelity} & \textbf{Modularity support} &  \textbf{Custom control} & \textbf{Use case for SRRS} \\
        \hline
        Gazebo+ ROS &
        High &
        Strong via URDF/SDF &
        ROS integration, plugins, Python/C++ nodes &
        Physical replication with mobile robots and manipulators \\
        \hline
        CoppeliaSim (V-REP) &
        High &
        Native modular modeling &
        Lua, Python /remote APIs &
        Detailed control of part interaction and modular assembly \\
        \hline
        PyBullet &
        Medium to High &
        Good (via URDF /SDF) &
        Python API (fast iteration) &
        Cube/voxel-based systems and contact-rich interactions \\
        \hline
        Webots &
        Medium &
        Moderate &
        C/C++ /Python controllers &
        Lightweight simulation of small-scale replication setups \\
        \hline
        Unity+ \ \ ML-Agents &
        Medium &
        High (custom prefabs) &
        C\#,Python with \ \ \ \ \ ML-agents &
        Learning-based or perception-driven replication strategies \\
        \hline
        Isaac Sim (NVIDIA) &
        High &
        Strong (via Omniverse) &
        Python scripting, ROS2 &
        High-fidelity simulation of complex autonomous robots \\
        \hline
        Golly  \ \ \ NetLogo  CA tools &
        None &
        Tile-based, symbolic &
        Rule-based/ scripting &
        Theoretical studies (e.g. von Neumann or Langton automata) \\
        \hline
    \end{tabular}
    \caption{Comparison of simulation environments for self-replicating robotic systems}
    \label{tab:tools}
\end{table}

In order to simulate kinematic self-replicating systems and to be able to implement
control functions that are also applicable to a physical system, the Gazebo simulator 
with ROS infrastructure seems to be an optimal choice.

\begin{figure}[h]
    \centering
    \includegraphics[width=15cm]{Gazebo.png} 
    \caption{Gazebo Harmonic simulation interface}
    \label{fig:Gazebo}
\end{figure}


Gazebo supports multiple physics engines, such as 
ODE, Bullet, and DART, allowing to model the dynamics of contact,
friction, collisions, and gravity with a high degree of accuracy. 

Modular robot design  can be modeled through ROS’s URDF (Unified Robot 
Description Format) and SDF (Simulation Description Format). These formats 
help to define robotic structures as combinations of interconnected 
parts and joints, closely mimicking the principles of modular self-replicating 
systems. Components can be described independently and reused across different 
robots or configurations. It support simulating a wide array of sensors, including
cameras, LiDAR, depth sensors, IMUs, proximity sensors, and contact detectors. 

Structured or unstructured environments for SRRS can be modeled in great detail,
placing objects with specific physical properties, and simulating real-world 
challenges like terrain irregularities or sensor noise.

ROS provides a powerful middleware layer that manages communication between 
different components of the robot and the environment. 
Through ROS, we can implement and test high-level control algorithms, including perception 
pipelines, decision-making logic, and task planning.
It supports distributed systems, so that multi-robot scenarios 
(such as cooperative replication) can be easily simulated. 
ROS also provides tools for recording and analyzing 
data, tuning controllers, and implementing feedback loops. 

Gazebo and ROS are open-source and highly extensible. Developers can create 
plugins for new sensors, actuators, or control strategies. This is particularly 
useful in self-replication research, where novel components or interfaces 
may need to be modeled. 
Additionally, the close integration between simulation 
and real-world systems means that code developed in Gazebo with ROS can 
often be deployed directly to real hardware with minimal changes. 


\section{Proposed approach}
Gazebo in combination with ROS can be used to create an interface between
simulated robot behavior with formal verification 
tools, such as NuSMV. Formal methods allow researchers to rigorously verify 
properties of robotic systems—such as correctness, safety, liveness, and reachability
—by constructing abstract models of their behaviors and validating them against 
logical specifications.

ROS-based systems define robotic structure and dynamics using URDF (Unified 
Robot Description Format) files and execute behavior through Python or C++ ROS 
nodes. 
These two sources—mechanical definitions and control logic—can be systematically 
translated into finite-state abstractions suitable for model checking with tools 
like NuSMV.

The process typically involves analyzing the state transitions implied by the 
ROS node logic, especially within action sequences (e.g., part pickup, alignment
, assembly), and discretizing them into states, inputs, and transitions. Each 
robot behavior or controller mode (e.g., "move", "grasp", "assemble") is 
represented as a finite state, while events and sensor conditions define transition 
rules. The mechanical configuration from URDF can define static properties or 
initial states in the system.

From this, a states diagram structure can be extracted, encoded 
in NuSMV's input language, and checked against temporal logic properties, described in 
CTL or LTL formulas. 
For instance, one could verify that a self-replicating 
robot always eventually reaches the assembly state (liveness) or that it never 
attempts to pick a module unless one is detected (safety).

Automating this translation—using Python scripts, that parse URDF and interpret ROS node 
behavior—is increasingly feasible, especially for systems designed with modular 
and rule-based architectures. Scripts or intermediate frameworks can analyze 
ROS bag files or logs to capture sequences of operations, and encode these as 
state transitions. 
% Some researchers also define behavior trees or state machines 
% in ROS using structured formats (e.g., SMACH, FlexBE), which can be directly 
% mapped to NuSMV models with minimal transformation.

This capability enables hybrid verification workflows, where simulation is used 
to test behavior under dynamic conditions, and formal verification is applied 
to prove correctness across all possible system executions. In the context of 
self-replication, this allows verification of critical properties like "a complete 
replica is only activated after all modules are assembled" or "robots do not 
collide while cooperating on part assembly".

By linking Gazebo/ROS to NuSMV, can be created a powerful toolchain to not 
only test and observe behavior in simulation but also prove correctness and 
robustness in all logical conditions—an essential feature for deploying self
-replicating systems in safety-critical or remote environments.

% \section{Analysis of results}


\chapter{Implementation}

\section{Implementation}

\subsection{Experimental setup}
As a subject of a research a simple robot will be taken, with two grippers 
and an elbow joint that can manipulate blocks (cuboctahedral voxels) one at a time.
An experimental platform will be realised as a simulation in Gazebo environment 
with control from ROS2 nodes. The  blocks and environment geometry robot joints 
and sensors and conections will be defined in Xacro/URDF files. 

    TODO: plan experiments: 
    Criteria: success rate, average time, \\
    1. Test with predefined assembly process (search, move, repeat) 
    with random distribution of the parts\\
    2. Test with predefined assembly process (...) with 
    , criteria and analysis 



% \subsubsection{Generation of a model and specifications for formal-verification}
% A custom program has to be written in Python, that parses the Xacro files and ROS2 Nodes controller to generate a NuSMV model:
%
  % - a Boolean occupancy matrix representing the neighbourhood on a field relevant to the current replication.
%
  % - replication program-counter plus local flags (gripper-occupied, block -held, etc.).
%
  % The liveness and safety properties has to be defined independent from the simulation, in a form close to the following:
  % \begin{itemize}
    % 
      % \item Liveness. $G (start \rightarrow F (Robot_1==ready \land Robot_2==ready))$ — once replication starts, eventually both robots are able to function independently.
%
      % \item Collision-free safety. $G \neg collision$ (or $AG \neg collision$ for CTL).
%
% \end{itemize}

